% =====================================================================
\part{Umfeld des Verfahrens}
% =====================================================================

\chapter{Codegenerierung und Build}
\label{chap:codegen}

Neben dem Lader erzeugt \texttt{sl} aus einem Metadatensatz die Schichten, die ein Projekt für die Nutzung auf mehreren Plattformen braucht: C-Header, C-Quelle, ein CMake-Modul, ein Python-Binding und die Build-Matrix (\texttt{codegen.py}), und es kompiliert und packt den Container (\texttt{build.py}).

\section{Ein durchgängiges Beispiel}

Die README zeigt die Metadaten der Bibliothek \texttt{wsat}:

\begin{lstlisting}[language=JSON, caption={Metadaten \texttt{wsat.json} (README, Schnellstart).}, label={lst:wsatjson}]
{
  "name": "wsat", "version": "1.2.0", "abi_version": 1,
  "targets": [{"platform": "linux", "arch": "x86_64"}],
  "dependencies": {},
  "exports": ["wsat_solve"],
  "signatures": {"wsat_solve": "int wsat_solve(const int *xs, unsigned int n)"},
  "init": "wsat_init", "fini": "wsat_fini"
}
\end{lstlisting}

Aus ihnen berechnet \texttt{generate\_header} den Header. Die folgende Ausgabe wurde durch Lesen des Generators abgeleitet, nicht durch einen Lauf erzeugt:

\begin{lstlisting}[language=C, caption={Abgeleitete Ausgabe von \texttt{generate\_header} für Listing~\ref{lst:wsatjson}.}, label={lst:header}]
/* Generated by sl for 'wsat' - do not edit. */
#ifndef WSAT_SL_H
#define WSAT_SL_H

#ifndef SL_EXPORT
#  if defined(_WIN32)
#    define SL_EXPORT __declspec(dllexport)
#  else
#    define SL_EXPORT __attribute__((visibility("default")))
#  endif
#endif

#ifdef __cplusplus
extern "C" {
#endif

#define WSAT_ABI_VERSION 1

SL_EXPORT int wsat_abi_version(void);
SL_EXPORT const char *wsat_sl_metadata(void);
SL_EXPORT int wsat_init(void);
SL_EXPORT int wsat_fini(void);
SL_EXPORT int wsat_solve(const int *xs, unsigned int n);

#ifdef __cplusplus
}
#endif

#endif /* WSAT_SL_H */
\end{lstlisting}

Das Python-Binding enthält für dieselben Metadaten die Tabelle \texttt{SIGNATURES = \{"wsat\_solve": ([ctypes.c\_void\_p, ctypes.c\_uint], ctypes.c\_int)\}}, denn \texttt{const int *} ist ein Zeiger, der nicht \texttt{char *} ist, und wird auf \texttt{c\_void\_p} abgebildet (Satz~\ref{satz:typen}).

\section{Eine Quelle der Wahrheit}

\begin{satz}[Konsistenz der erzeugten Schichten]
	\label{satz:schichtkonsistenz}
	Jeder Generator ist eine reine Funktion des Metadatensatzes $M$: Er liest keine anderen Eingaben und erzeugt für gleiche $M$ gleiche Ausgaben. Der Header deklariert genau die Menge
	\[
	\{\nu\texttt{\_abi\_version},\ \nu\texttt{\_sl\_metadata}\}\cup\{\iota,\varphi\}\setminus\{\texttt{None}\}\cup X,
	\]
	das Python-Binding enthält einen Eintrag je Element von $X$ in der Reihenfolge der Exportliste, und die Build-Matrix hat genau $|T|$ Einträge, einen je Target in der Ordnung von $T$.
\end{satz}

\begin{proof}
	Die Funktionen \texttt{generate\_header}, \texttt{generate\_c\_source}, \texttt{generate\_cmake}, \texttt{generate\_python\_binding} und \texttt{build\_matrix} haben als einzigen nichtkonstanten Parameter \texttt{meta} (und \texttt{compiler} mit Vorgabe \texttt{"cc"}); sie verwenden \texttt{json.dumps(\ldots, sort\_keys=True)} und Schleifen über \texttt{meta.exports} und \texttt{meta.targets}, deren Reihenfolge feststeht. Die Menge der Deklarationen ist durch die Zeilen in \texttt{generate\_header} gegeben: zwei feste Symbole, je ein Hook, falls nicht \texttt{None}, und eine Zeile je Export. Der Binding-Generator iteriert über \texttt{meta.exports} und baut je Symbol eine Zeile; die Matrix ist eine Listenkomprehension über \texttt{metadata.targets}.
\end{proof}

\begin{lemma}[Eingebettete Metadaten]
	\label{lem:eingebettet}
	Die vom C-Quelltext zurückgegebene Zeichenkette \texttt{<name>\_sl\_metadata()} ist ein gültiges C-Zeichenkettenliteral und gleich \texttt{json.dumps(M.to\_dict(), sort\_keys=True)}. Sie ist byteweise gleich der Metadaten-Kodierung im Container genau dann, wenn die Targettupel von $M$ und der beim Packen verwendeten Metadaten gleich sind.
\end{lemma}

\begin{proof}
	\texttt{json.dumps} mit der Vorgabe \texttt{ensure\_ascii=True} erzeugt nur druckbare ASCII-Zeichen; Steuerzeichen, Zeilenumbrüche und Nicht-ASCII-Zeichen werden als \texttt{\textbackslash n}, \texttt{\textbackslash uXXXX} usw. geschrieben. \texttt{\_c\_string} verdoppelt Rückstriche und maskiert Anführungszeichen; das Ergebnis ist ein gültiges Literal, das als Zeichenkette den ursprünglichen Text ergibt. Der Container verwendet denselben Ausdruck (\texttt{pack}), jedoch mit den Metadaten, die an \texttt{pack} übergeben werden; \texttt{pack\_files} übergibt \texttt{metadata.with\_targets(payloads)}, dessen Targets aufsteigend sortiert sind und der Menge der Payloads entsprechen. Alle anderen Felder stimmen überein; ein Unterschied kann deshalb nur im Feld \texttt{targets} auftreten, und er tritt auf, wenn die Tupel verschieden sind (andere Reihenfolge oder andere Menge).
\end{proof}

Der Lader liest die eingebettete Zeichenkette nicht; sie dient der Selbstauskunft des Payloads. Die Folge ist, dass das Feld \texttt{targets} der Metadaten bei \texttt{sl pack} und \texttt{sl build} nur beratend ist: Das Targetfeld des Containers wird aus den tatsächlich vorhandenen Payloads gebildet.

\section{Signaturen und Typabbildung}

\begin{definition}[Signaturfragment]
	\label{def:fragment}
	Das \emph{Fragment} $\mathcal{F}_{\mathrm{sig}}$ besteht aus Zeichenketten der Form $R\ n\,(\,a_1,\ldots,a_k\,)\,[;]$ oder $R\ n\,(\texttt{void})\,[;]$, wobei $R$ und die $a_i$ jeweils aus einem Basistyp der Tabelle $\mathcal{T}$ (22 Einträge: \texttt{void}, \texttt{char}, \texttt{short}, \texttt{int}, \texttt{unsigned}, \texttt{unsigned int}, \texttt{long}, \texttt{unsigned long}, \texttt{long long}, \texttt{unsigned long long}, \texttt{float}, \texttt{double}, \texttt{size\_t}, \texttt{bool} und die acht Typen \texttt{int8\_t} bis \texttt{uint64\_t}), optionalem \texttt{const} und beliebig vielen \texttt{*} bestehen, jedes $a_i$ zusätzlich aus einem optionalen Parameternamen, der ein Bezeichner und kein Typwort ist; Klammern kommen sonst nicht vor.
\end{definition}

\begin{lemma}[Korrektheit von \texttt{parse\_signature}]
	\label{lem:parse}
	Für $s\in\mathcal{F}_{\mathrm{sig}}$ mit Funktionsname $n$ liefert \texttt{parse\_signature}$(s,n)$ den Rückgabetyp und die Liste der Argumenttypen ohne Parameternamen, in der Reihenfolge der Deklaration; für $(\texttt{void})$ ist die Liste leer. Für $s$ mit anderem Namen wird \texttt{ValueError} geworfen.
\end{lemma}

\begin{proof}
	Der reguläre Ausdruck \texttt{\^{}(?P<ret>.*?[\textbackslash s*])(?P<n>[A-Za-z\_]\textbackslash w*)\textbackslash s*\textbackslash((?P<args>.*)\textbackslash)\textbackslash s*;?\textbackslash s*\$} verlangt, dass der Name unmittelbar vor der ersten öffnenden Klammer steht. In $\mathcal{F}_{\mathrm{sig}}$ kommt vor ihr keine Klammer vor, also ist \texttt{ret} der Text bis vor den letzten Bezeichner und der Name der Bezeichner; \texttt{args} reicht bis zur einzigen schließenden Klammer. Ein abweichender Name löst die Prüfung \texttt{match.group("n") != symbol} aus. Die Argumente werden an Kommata getrennt, was ohne Klammern und ohne Kommata in Typen eindeutig ist. In jedem Argument trennt der Code \texttt{*} durch Leerzeichen ab und zerlegt in Wörter. Das letzte Wort wird als Parametername entfernt, wenn es nicht \texttt{*} und kein Typwort ist und mehr als ein Wort vorhanden ist. Für einen Typ aus $\mathcal{T}$ ohne Namen ist das letzte Wort ein Typwort oder ein \texttt{*}, es wird nicht entfernt; mit Namen ist es der Name (kein Typwort nach Voraussetzung) und wird entfernt. Ein einzelnes Wort wird nie entfernt. Ist die Argumentliste genau \texttt{void}, entsteht die leere Liste.
\end{proof}

\begin{annahme}[Größen in \texttt{ctypes}]
	\label{ann:ctypes}
	Für jedes $\tau\in\mathcal{T}$ stimmt der \texttt{ctypes}-Typ $\kappa(\tau)$ in Größe, Vorzeichenbehaftung und Argumentübergabe mit dem C-Typ $\tau$ der Zielplattform überein (so definiert die \texttt{ctypes}-Dokumentation etwa \texttt{c\_long} als C-\texttt{long} der Plattform), und Zeiger haben die Größe von \texttt{void *}.
\end{annahme}

\begin{satz}[Typabbildung]
	\label{satz:typen}
	Unter Annahme~\ref{ann:ctypes} gilt für \texttt{ctype\_expr}: Ein Typ mit mindestens einem \texttt{*} wird, nach Entfernen von \texttt{const}, genau dann auf \texttt{c\_char\_p} abgebildet, wenn seine Wörter genau \texttt{char}, \texttt{*} sind; sonst auf \texttt{c\_void\_p}. Ein Typ ohne \texttt{*} wird auf den Eintrag der Tabelle abgebildet, wenn er in $\mathcal{T}$ liegt, sonst wird \texttt{ValueError} geworfen. Insbesondere werden \texttt{unsigned char}, \texttt{unsigned short}, \texttt{long int} und \texttt{long double} abgelehnt, und die Plattformabhängigkeit von \texttt{long} (32 Bit unter Windows, 64 Bit unter LP64-Systemen) wird von \texttt{ctypes} aufgefangen, nicht vom Generator.
\end{satz}

\begin{proof}
	Der erste Teil ist die Fallunterscheidung im Code (nach \texttt{if "*" in tokens}). Der zweite Teil ist die Suche in \texttt{\_SCALARS}; die genannten Typen stehen nicht in der Tabelle. Dass \texttt{c\_long} die Plattformgröße trägt, ist Annahme~\ref{ann:ctypes}; der Generator verwendet für \texttt{long} das Symbol, nicht eine feste Breite.
\end{proof}

\begin{bemerkung}
	Die Abbildung \emph{jedes} \texttt{char *}-Rückgabewerts auf \texttt{c\_char\_p} macht aus dem Zeiger eine Python-Bytefolge (Kopie bis zum Nullbyte); der Zeigerwert geht verloren. Gibt die Funktion Speicher mit Eigentum an den Aufrufer zurück, kann er nicht freigegeben werden. Der Entwurf nennt dazu die Funktion \texttt{wsat\_free}; eine Notation für Besitzverhältnisse in den Signaturen fehlt (Ausblick).
\end{bemerkung}

\section{Build-Orchestrierung}

\begin{satz}[Ergebnis des Builds]
	\label{satz:build}
	\texttt{build\_sl}$(M,\mathrm{sources},o,\ldots,\mathrm{extra})$ erzeugt einen Container, dessen Targetmenge gleich $\mathrm{dom}(\mathrm{extra})\cup\{h\}$ ist, wobei $h$ das Target des Hosts ist; der Payload zu $h$ ist die frisch kompilierte Bibliothek, auch wenn \texttt{extra} einen Eintrag für $h$ enthält. Das Feld \texttt{targets} von $M$ geht nicht in das Ergebnis ein.
\end{satz}

\begin{proof}
	In \texttt{build\_sl} wird \texttt{files = \{**extra, host: lib\}} gebildet; bei gleichem Schlüssel gewinnt der spätere Eintrag, also \texttt{lib}. \texttt{pack\_files} ruft \texttt{pack(metadata.with\_targets(payloads), payloads)}, ersetzt also die Targetliste der Metadaten durch die Schlüssel von \texttt{payloads}.
\end{proof}

Zwei Einschränkungen des Builds sind für die Beispielprojekte erheblich. Erstens verwendet \texttt{toolchain\_flags} für Windows nur \texttt{-shared}, also die Form von GCC und Clang; \texttt{find\_compiler} nennt aber auch \texttt{cl} als Kandidaten, dessen Aufrufform anders ist (\texttt{/LD}, \texttt{/Fe}). Zweitens können Compiler, Sprachstandard und Bibliotheken nicht angegeben werden: \texttt{-I} und \texttt{-D} sind die einzigen Zusätze, und \texttt{cc} linkt keine C++-Laufzeit. Ein C++17-Projekt wie \repo{csubgraph} lässt sich damit nicht ohne Umweg bauen (Befund~\ref{bef:build}). Das ist keine Schwäche des Formats, sondern eine Lücke des Werkzeugs, die der Ausblick schließt.

\chapter{Versionen: Konvexität und Komplexität}
\label{chap:versionen}

Die Metadaten tragen Versionsbereiche für Abhängigkeiten. Dieses Kapitel beantwortet zwei Fragen: Welche Mengen von Versionen lassen sich ausdrücken? Und warum wählt der Lader nicht unter mehreren Kandidaten aus?

\section{Die Bereichssprache}

\begin{satz}[Ordnung und Konvexität]
	\label{satz:konvex}
	\begin{enumerate}[label=(\alph*)]
		\item $(\N^3,<)$ mit der lexikographischen Ordnung ist total, und \texttt{satisfies}$(v,\mathsf{s})$ wird in Zeit $O(|\mathsf{s}|)$ entschieden.
		\item Für jeden Bereich $\mathsf{s}$ ist $S(\mathsf{s})$ \emph{konvex}: Aus $u<v<w$ und $u,w\in S(\mathsf{s})$ folgt $v\in S(\mathsf{s})$.
		\item Die Bereichssprache ist unter Konjunktion abgeschlossen (Aneinanderreihen mit Komma), nicht aber unter Disjunktion: Es gibt keinen Bereich $\mathsf{s}$ mit $S(\mathsf{s})=S(\texttt{<1.0.0})\cup S(\texttt{>=2.0.0})$.
		\item Es gibt unerfüllbare Bereiche mit $S(\mathsf{s})=\emptyset$, etwa \texttt{>=2.0.0,<1.0.0}, und \texttt{Metadata} akzeptiert sie.
	\end{enumerate}
\end{satz}

\begin{proof}
	(a) Die lexikographische Ordnung auf $\N^3$ ist total: Für $v\ne w$ gibt es eine erste Komponente mit Unterschied, und dort ist eine kleiner. Die Prüfung vergleicht $v$ mit jeder Schranke einmal, jeder Vergleich ist $O(1)$.
	(b) Jede einzelne Einschränkung hat eine konvexe Lösungsmenge: $\{v\ge b\}$ und $\{v>b\}$ enthalten mit $u$ auch alles Größere, $\{v\le b\}$ und $\{v<b\}$ mit $w$ auch alles Kleinere, $\{v=b\}$ ist einelementig; in jedem Fall liegt $v$ zwischen zwei Elementen in der Menge. Der Durchschnitt endlich vieler konvexer Mengen ist konvex: Liegen $u,w$ im Durchschnitt, so liegt $v$ in jeder der Mengen.
	(c) Konjunktion: $S(\mathsf{s}_1,\mathsf{s}_2)=S(\mathsf{s}_1)\cap S(\mathsf{s}_2)$ nach Definition. Disjunktion: $u=(0,5,0)$ und $w=(2,0,0)$ liegen in der Vereinigung, $v=(1,0,0)$ liegt zwischen ihnen und nicht darin; nach (b) kann die Vereinigung kein $S(\mathsf{s})$ sein.
	(d) Für $v\ge(2,0,0)$ und $v<(1,0,0)$ gibt es kein $v$. Die Prüfung in \texttt{Metadata.\_\_post\_init\_\_} ruft nur \texttt{\_parse\_range} auf, das Syntax, nicht Erfüllbarkeit prüft.
\end{proof}

Zwei Folgerungen: Eine Abhängigkeit mit unerfüllbarem Bereich führt bei jedem \texttt{open} zu \texttt{DEPENDENCY\_MISSING}; eine Prüfung beim Erzeugen der Metadaten wäre billig (Ausblick). Und der Entwurf-Beispielbereich \texttt{>=1.2,<2} (Anhang~A) ist in der Implementierung syntaktisch ungültig, weil sie dreiteilige Versionen verlangt (Befund~\ref{bef:abw}).

\section{Auswahl unter mehreren Kandidaten ist NP-vollständig}

Der Lader nimmt den ersten Treffer (Satz~\ref{satz:auswahl}) und prüft nur ihn. Man könnte meinen, das sei eine Bequemlichkeit. Der folgende Satz zeigt, dass es eine Notwendigkeit ist: Würde der Lader unter \emph{allen} gefundenen Dateien eine gemeinsame, verträgliche Auswahl suchen, wäre das Problem NP-vollständig.

\begin{definition}[Auswahlproblem \textsc{DepSelect}]
	\label{def:depselect}
	Eine Instanz besteht aus einer endlichen Menge von Paketen $\Pi$; jedem Paket $\pi$ ist eine endliche Menge $V(\pi)\subset\N^3$ von Versionen zugeordnet; jede Version $(\pi,v)$ trägt eine endliche Liste von Anforderungen $(\pi',\mathsf{s})$ mit einem Bereich $\mathsf{s}$; und es gibt eine Wurzel mit einer Liste von Anforderungen. Eine \emph{Auswahl} ist eine partielle Abbildung $c$ von Paketen auf Versionen, die jedem ausgewählten Paket höchstens eine Version zuordnet. Sie ist \emph{zulässig}, wenn für jede Anforderung $(\pi',\mathsf{s})$ der Wurzel und jeder ausgewählten Version $(\pi,c(\pi))$ gilt: $\pi'$ ist ausgewählt und $c(\pi')\models\mathsf{s}$. Die Frage lautet, ob eine zulässige Auswahl existiert.
\end{definition}

\begin{satz}[\textsc{DepSelect} ist NP-vollständig]
	\label{satz:npvoll}
	\textsc{DepSelect} ist NP-vollständig, auch wenn jedes Paket höchstens drei Versionen und jede Version höchstens eine Anforderung hat und alle Bereiche die Form \texttt{>=2.0.0}, \texttt{<2.0.0} oder \texttt{*} haben.
\end{satz}

\begin{proof}
	\emph{In NP.} Eine Auswahl hat polynomielle Größe; ihre Zulässigkeit prüft man mit einem Durchlauf über die Anforderungen und je einer Bereichsprüfung (Satz~\ref{satz:konvex}(a)).

	\emph{NP-schwer.} Wir reduzieren 3-SAT (Karp \cite{karp1972}). Sei $\varphi=C_1\wedge\cdots\wedge C_m$ über $x_1,\ldots,x_n$, jede Klausel mit genau drei Literalen (kürzere Klauseln durch Wiederholen eines Literals aufgefüllt). Wir bilden folgende Instanz.
	\begin{itemize}
		\item Pakete $X_1,\ldots,X_n$ mit den Versionen $1.0.0$ („falsch“) und $2.0.0$ („wahr“), ohne Anforderungen.
		\item Pakete $K_1,\ldots,K_m$ mit den Versionen $1.0.0$, $2.0.0$, $3.0.0$. Ist das $t$-te Literal von $C_j$ gleich $x_i$, so fordert die Version $t.0.0$ von $K_j$ das Paket $X_i$ mit \texttt{>=2.0.0}; ist es $\neg x_i$, so fordert sie $X_i$ mit \texttt{<2.0.0}.
		\item Die Wurzel fordert jedes $K_j$ mit \texttt{*}.
	\end{itemize}
	Die Größe ist linear in $|\varphi|$.

	($\Rightarrow$) Sei $\alpha$ eine erfüllende Belegung. Wähle für jedes $j$ ein $t_j$, dessen Literal in $C_j$ unter $\alpha$ wahr ist, und setze $c(K_j)=t_j.0.0$. Setze $c(X_i)=2.0.0$, falls $\alpha(x_i)=\mathrm{wahr}$, sonst $1.0.0$, für alle $X_i$, die von einer gewählten Version gefordert werden. Die Anforderungen der Wurzel sind erfüllt; die Anforderung der gewählten Version von $K_j$ ist es, weil das Literal wahr ist: bei $x_i$ ist $c(X_i)=2.0.0\ge2.0.0$, bei $\neg x_i$ ist $c(X_i)=1.0.0<2.0.0$. Die $X_i$ haben keine Anforderungen. Also ist $c$ zulässig.

	($\Leftarrow$) Sei $c$ zulässig. Setze $\alpha(x_i)=\mathrm{wahr}$ genau dann, wenn $X_i$ ausgewählt ist und $c(X_i)=2.0.0$. Für jedes $j$ ist $K_j$ ausgewählt (Wurzel), etwa mit $c(K_j)=t.0.0$. Dessen Anforderung verlangt für das $t$-te Literal $\ell$ von $C_j$: $X_i$ ist ausgewählt und $c(X_i)\ge2.0.0$ (bei $\ell=x_i$), also $\alpha(x_i)=\mathrm{wahr}$ und $\ell$ wahr; oder $c(X_i)<2.0.0$ (bei $\ell=\neg x_i$), also $\alpha(x_i)=\mathrm{falsch}$ und $\ell$ wahr. Jede Klausel hat ein wahres Literal, also erfüllt $\alpha$ die Formel $\varphi$.
\end{proof}

\begin{korollar}[Eindeutige Auflösung ist notwendig für Polynomialzeit]
	\label{kor:eindeutig}
	Gäbe es einen Lader, der bei mehreren Kandidatendateien je Name stets eine zulässige Auswahl findet, wenn eine existiert, so entschiede er \textsc{DepSelect}; er läuft dann im ungünstigsten Fall nicht in polynomieller Zeit, sofern $\mathrm{P}\ne\mathrm{NP}$.
\end{korollar}

\begin{proof}
	Aus Satz~\ref{satz:npvoll} und der Definition der NP-Vollständigkeit.
\end{proof}

\begin{satz}[Unvollständigkeit, Polynomialität des Laders]
	\label{satz:unvollstaendig}
	Der Lader löst jeden Namen durch \texttt{locate} zu einer einzigen Datei auf (Satz~\ref{satz:auswahl}).
	\begin{enumerate}[label=(\alph*)]
		\item Es gibt eine Konfiguration, in der eine Datei, die einen Bereich erfüllt, im Suchpfad liegt, und \texttt{open} dennoch mit \texttt{DEPENDENCY\_MISSING} scheitert.
		\item \texttt{open}$(r)$ läuft in Zeit $O\bigl((1+|E_r|)\cdot n_{\max}+\sum_{f\in V_r}|f|\bigr)$ mit den Dateigrößen $|f|$ und den Kosten der Bereichsprüfungen und Hooks nicht gerechnet.
	\end{enumerate}
\end{satz}

\begin{proof}
	(a) Sei \texttt{SL\_PATH}$=d_1\texttt{:}d_2$; \texttt{d1/base.sl} habe die Version $1.0.0$, \texttt{d2/base.sl} die Version $2.0.0$, und \texttt{top} fordere \texttt{base} mit \texttt{>=2.0.0}. \texttt{locate} liefert die Datei in $d_1$; die Bereichsprüfung scheitert und \texttt{\_open\_dependency} lässt den Fehler weiterlaufen. Die Datei in $d_2$ wird nie betrachtet. (b) Folgt aus Satz~\ref{satz:kosten} und der Tatsache, dass \texttt{\_read} und \texttt{unpack} jede Datei einmal in linearer Zeit lesen und zerlegen.
\end{proof}

Der Zusammenhang zu den übrigen Arbeiten der Reihe ist keine Zufälligkeit: Die Kette 3-SAT $\to$ Clique $\to$ Teilgraph der Vorarbeit und die Reduktion in Satz~\ref{satz:npvoll} sind dieselbe Art von Aussage über verschiedene Gegenstände. Dass der Entwurf bei widersprüchlichen Kandidaten einen Fehler verlangt, statt zu suchen, ist folglich auch aus Sicht der Berechenbarkeit gut begründet; die Frage, wie man die Diagnose bei Mehrdeutigkeit verbessert, ohne zu suchen, behandelt der Ausblick.

\chapter{Sicherheit}
\label{chap:sicherheit}

\section{Angreifermodell}

Betrachtet werden drei Angreiferfähigkeiten: (A1) Der Angreifer kann Umgebungsvariablen des Prozesses setzen (nicht aber das Programm ändern). (A2) Er kann Dateien in Verzeichnissen anlegen, die für alle beschreibbar sind. (A3) Er kann Dateien zwischen zwei Zugriffen des Laders ersetzen, soweit seine Rechte es zulassen. Ziel ist, den Prozess zum Laden von Code zu bringen, den der Betreiber nicht vorgesehen hat.

\section{Die klassischen Angriffswege}

\begin{table}[!ht]
	\centering
	\small
	\begin{tabularx}{\textwidth}{@{}L{3.6cm}L{4.4cm}Y@{}}
		\toprule
		Angriffsweg & Maßnahme in \texttt{sl} & Bleibender Rest \\
		\midrule
		Arbeitsverzeichnis (DLL-Hijacking) & Satz~\ref{satz:cwd}: nicht durchsucht, wenn $R$, $A$, $\Sigma_{\mathrm{sys}}$ absolut & $R$ und $A$ werden nicht normalisiert (B2) \\[0.3em]
		\texttt{PATH}, \texttt{LD\_LIBRARY\_PATH}, \texttt{DYLD\_LIBRARY\_PATH} & Satz~\ref{satz:struktur}: nie gelesen & Windows liest \texttt{ProgramData}, \texttt{ProgramFiles} für $\Sigma_{\mathrm{sys}}$ \\[0.3em]
		\texttt{SL\_PATH} (A1) & Satz~\ref{satz:gehaertet}(i): im gehärteten Modus ignoriert & im Normalbetrieb vertraut der Lader der Variable \\[0.3em]
		Beschreibbare Verzeichnisse (A2) & Satz~\ref{satz:gehaertet}(ii), (iii): Ablehnung, fail closed & nur unmittelbares Verzeichnis, nur \emph{others}, Fehler bei \texttt{stat} lässt zu (B3) \\[0.3em]
		Veränderter Payload & Satz~\ref{satz:integritaet}: Hash & kein Schutz vor gezielter Veränderung (Bemerkung~\ref{bem:hashgrenze}) \\[0.3em]
		Nicht deklarierte Symbole & Satz~\ref{satz:export} & keine Sicherheitsgrenze im Prozess \\
		\bottomrule
	\end{tabularx}
	\caption{Angriffswege, Gegenmaßnahmen und Rest.}
	\label{tab:angriff}
\end{table}

\section{Einmaliges Lesen}
\label{sec:toctou}

\begin{lemma}[Validierte Bytes sind die abgebildeten Bytes]
	\label{lem:einmal}
	Die Bytes, die \texttt{unpack} prüft, sind dieselben, aus denen der Payload geschnitten, auf Hash geprüft und in die temporäre Datei geschrieben wird.
\end{lemma}

\begin{proof}
	\texttt{Loader.\_load} liest die Datei einmal (\texttt{self.\_read(path)}) und übergibt die Byte-Folge an \texttt{unpack}; \texttt{payload\_for} schneidet aus ihrem Payload-Bereich; \texttt{map\_payload} schreibt genau diese Folge. Es gibt keinen zweiten Lesezugriff auf den Pfad.
\end{proof}

Das Lemma schließt eine häufige Klasse von Wettläufen aus (Prüfen eines Pfades, späteres Öffnen eines anderen Inhalts). Offen bleibt der Wettlauf zwischen der Verzeichnisprüfung in \texttt{locate} und dem Lesen in \texttt{\_read}, und das Öffnen der temporären Datei durch den Systemlader nach dem Schreiben. Unter POSIX ist das Verzeichnis nach der Dokumentation von \texttt{mkdtemp} nur für den Besitzer zugänglich; unter Windows hängt der Schutz von den Zugriffsrechten des übergeordneten Verzeichnisses ab \offen{nicht geprüft}. Da der Ort über \texttt{TMPDIR}, \texttt{TEMP} oder \texttt{TMP} bestimmt wird, gehört auch diese Variable zum Vertrauensbereich des Prozesses (A1).

\chapter{Anwendung auf die Beispielprojekte}
\label{chap:anwendung}

\section{Vorher und nachher}

Der Entwurf zeigt die Wirkung an einem Beispiel (Abschnitt~12):

\begin{lstlisting}[language=Python, caption={Vorher: plattformabhängiger Dateiname (nach \texttt{proposal.md}, Abschnitt~12).}, label={lst:vorher}]
def library_filename() -> str:
    if sys.platform == "win32":
        return "wsat_shim.dll"
    if sys.platform == "darwin":
        return "libwsat_shim.dylib"
    return "libwsat_shim.so"
\end{lstlisting}

\begin{lstlisting}[language=Python, caption={Nachher: dieselbe Zeile auf jedem System.}, label={lst:nachher}]
import sl
lib = sl.open("wsat_shim")      # same code on every system
\end{lstlisting}

Tabelle~\ref{tab:ersatz} prüft, welche der sechs Stellen aus Abschnitt~\ref{sec:evidenz} das Paket \texttt{sl} \emph{tatsächlich} ersetzt.

\begin{table}[!ht]
	\centering
	\small
	\begin{tabularx}{\textwidth}{@{}L{4.3cm}L{3.3cm}Y@{}}
		\toprule
		Stelle & Ersetzt? & Begründung \\
		\midrule
		\texttt{library\_filename()} mit vier Namen & ja & logischer Name, Lemma~\ref{lem:name}; Voraussetzung: die Bibliothek wird als \SL\ gebaut \\[0.3em]
		\texttt{add\_dll\_directory} für \texttt{PATH}-Einträge & für deklarierte Abhängigkeiten ja & Abhängigkeiten werden nach derselben Ordnung aufgelöst (Satz~\ref{satz:struktur}); native Laufzeit-DLLs sind nicht erfasst \\[0.3em]
		MinGW-Sonderfall in \texttt{build\_shim()} & teilweise & \texttt{sl build} kennt nur GCC/Clang-Flags (Satz~\ref{satz:build}, B9) \\[0.3em]
		Makro \texttt{CSUB\_EXPORT} & ja & erzeugter Header mit \texttt{SL\_EXPORT} (Listing~\ref{lst:header}) \\[0.3em]
		statisches Linken der MinGW-Laufzeit & nein & das Feld \texttt{runtime} hat keine Wirkung (B11) \\[0.3em]
		Plattformteile der README & teilweise & \texttt{sl build} ist überall gleich; Umgebungsvariablen und Shell bleiben verschieden \\
		\bottomrule
	\end{tabularx}
	\caption{Ersatz der sechs Plattformstellen durch \texttt{sl} (Befunde in Kapitel~\ref{chap:befunde}).}
	\label{tab:ersatz}
\end{table}

Die Tabelle ist ernüchternder als der Entwurf: Von sechs Stellen entfallen zwei ganz, drei teilweise, eine gar nicht. Das ist ein Ergebnis dieser Arbeit und kein Einwand gegen den Entwurf: Es zeigt, welche Teile des Referenzpakets noch ausgebaut werden müssen, damit die Zusage des Entwurfs trägt (Ausblick, Abschnitte~\ref{sec:runtime-ausblick} und~\ref{sec:migration}).

\section{Metadaten für \texttt{wsat\_shim}}

Der Entwurf skizziert die Metadaten der Bibliothek in TOML (Anhang~A); in der Implementierung sind sie ein JSON-Objekt. Die Entsprechung:

\begin{lstlisting}[language=TOML, caption={Metadatenskizze (\texttt{proposal.md}, Anhang~A, gekürzt).}, label={lst:toml}]
[sl]
container_version = 1
name        = "wsat_shim"
version     = "1.4.0"
abi_version = 3

[[target]]
os = "linux"
arch = "x86_64"
payload = "elf"

[runtime]
cpp = "embedded"

[[dependency]]
name = "cryptosolver"
version = ">=1.2,<2"

[exports]
symbols = ["wsat_abi_version", "wsat_solve", "wsat_free"]

[lifecycle]
init = "wsat_init"
fini = "wsat_fini"
\end{lstlisting}

Die Unterschiede zur Implementierung: \texttt{payload} (\texttt{elf}) ist dort nicht Teil der Metadaten, weil sich das Format aus dem Target ergibt; \texttt{cpp} heißt \texttt{runtime} mit den Werten \texttt{embedded} und \texttt{external}; Abhängigkeiten sind ein Wörterbuch Name $\to$ Bereich; der Bereich \texttt{>=1.2,<2} müsste \texttt{>=1.2.0,<2.0.0} heißen; \texttt{container\_version} steht im Kopf und nicht in den Metadaten.

\section{Ablauf eines \texttt{open} mit Abhängigkeit}

Abbildung~\ref{abb:sequenz} zeigt den Ablauf von \texttt{sl.open("top")} für eine Bibliothek \texttt{top}, die \texttt{base} deklariert. Er entspricht dem Test \texttt{test\_dependency\_chain\_and\_release} (dort mit einer Zwischenbibliothek mehr).

\begin{figure}[!ht]
	\centering
	\resizebox{0.98\textwidth}{!}{%
	\begin{tikzpicture}[yscale=1.3, font=\footnotesize, >={Stealth[length=2mm]},
		head/.style={draw=mainblue, fill=lightgray, minimum width=2.1cm, minimum height=0.7cm, align=center},
		life/.style={dashed, gray},
		msg/.style={->, thick},
		ret/.style={->, dashed}]
		\def\xa{0} \def\xb{3} \def\xc{6} \def\xd{9} \def\xe{12}
		\node[head] (a) at (\xa,0) {Anwendung};
		\node[head] (b) at (\xb,0) {Loader};
		\node[head] (c) at (\xc,0) {search / container};
		\node[head] (d) at (\xd,0) {Adapter};
		\node[head] (e) at (\xe,0) {Hooks (nativ)};
		\foreach \x in {\xa,\xb,\xc,\xd,\xe} \draw[life] (\x,-0.4) -- (\x,-13.2);
		\draw[msg] (\xa,-1) -- (\xb,-1) node[midway, above] {open("top")};
		\draw[msg] (\xb,-1.7) -- (\xc,-1.7) node[midway, above] {locate("top")};
		\draw[ret] (\xc,-2.3) -- (\xb,-2.3) node[midway, above] {Pfad};
		\draw[msg] (\xb,-3.0) -- (\xc,-3.0) node[midway, above] {unpack(read(Pfad))};
		\draw[ret] (\xc,-3.6) -- (\xb,-3.6) node[midway, above] {Metadaten, Payload};
		\node[draw=mainblue, fill=white, align=left, font=\scriptsize] at (\xb,-4.35) {ABI-Bereich, Target, Versionsbereich prüfen};
		\draw[msg] (\xb,-5.1) -- (\xc,-5.1) node[midway, above] {locate("base", dir(top))};
		\draw[ret] (\xc,-5.7) -- (\xb,-5.7) node[midway, above] {Pfad};
		\node[draw=mainblue, fill=white, align=left, font=\scriptsize] at (\xb,-6.4) {rekursiv: \texttt{\_load(base)} wie oben};
		\draw[msg] (\xb,-7.1) -- (\xd,-7.1) node[midway, above] {map\_payload(base)};
		\draw[msg] (\xb,-7.8) -- (\xe,-7.8) node[pos=0.5, above] {base\_init()};
		\draw[ret] (\xe,-8.3) -- (\xb,-8.3) node[midway, above] {0};
		\draw[msg] (\xb,-9.0) -- (\xd,-9.0) node[midway, above] {map\_payload(top)};
		\draw[msg] (\xb,-9.7) -- (\xe,-9.7) node[pos=0.5, above] {top\_init()};
		\draw[ret] (\xe,-10.2) -- (\xb,-10.2) node[midway, above] {0};
		\node[draw=mainblue, fill=white, align=left, font=\scriptsize] at (\xb,-10.9) {$\Lambda\mathrel{+}=\{\text{base},\text{top}\}$, rc(top)=1, rc(base)=1};
		\draw[ret] (\xb,-11.7) -- (\xa,-11.7) node[midway, above] {Handle(top)};
		\node[font=\scriptsize, text=darkgray, align=left, anchor=west] at (\xa,-12.7) {close(top): fini(top), unmap(top), dann fini(base), unmap(base)};
	\end{tikzpicture}}
	\caption{Ablauf von \texttt{open("top")} mit der Abhängigkeit \texttt{base}. Init-Reihenfolge nach Satz~\ref{satz:reihenfolge}(i); die Freigabe läuft in umgekehrter Reihenfolge (Satz~\ref{satz:reihenfolge}(iii)).}
	\label{abb:sequenz}
\end{figure}

\section{Der Test als Beleg: Reihenfolge der Abbildungen}

Der Test \texttt{test\_dependency\_chain\_and\_release} (\texttt{tests/test\_loader.py}) prüft für die Kette \texttt{top} $\to$ \texttt{mid} $\to$ \texttt{base} mit dem Fake-Adapter, dass \texttt{map} in der Reihenfolge \texttt{base}, \texttt{mid}, \texttt{top} und \texttt{unmap} in der Reihenfolge \texttt{top}, \texttt{mid}, \texttt{base} geschieht:

\begin{lstlisting}[language=Python, caption={Ausschnitt aus \texttt{tests/test\_loader.py} (Zeilen 92--100).}, label={lst:test}]
def test_dependency_chain_and_release(loader, adapter, write_sl):
    write_sl("base")
    write_sl("mid", dependencies={"base": ">=1.0.0"})
    write_sl("top", dependencies={"mid": "*"})
    handle = loader.open("top")
    assert adapter.mapped == ["base", "mid", "top"]       # Abhaengigkeiten zuerst
    assert loader.info(handle)["dependencies"] == ["mid"]
    loader.close(handle)
    assert adapter.unmapped == ["top", "mid", "base"]
\end{lstlisting}

Das ist ein Einzelfall der Aussagen (i) und (iii) von Satz~\ref{satz:reihenfolge}; der Satz gilt für jeden azyklischen Graphen. Der Test belegt also eine Instanz, der Satz die Regel.

